\section{Extended StampCheck Setup and Costs}
\label{app:stamp-based}
\paragraph{Facts and mutations.} Table~\ref{tab:facts} lists 16 base facts and 3 derived facts, each standing in for a critical use (text a contact, reply to the newest message, call back, share the newest photo, join a meeting, act on a setting, edit a note).

\begin{table}[t]
\centering
\caption{Remembered facts and injected mutations. agg.\ = aggregate. Derived facts: Dana+stand-up, newest photo+Dana, newest SMS+Omar.}
\label{tab:facts}
\mtab\scriptsize
\setlength{\tabcolsep}{2.5pt}
\begin{tabular}{p{1.25cm}p{3.75cm}p{2.85cm}}
\toprule
\mhead Source & Facts & Mutations (another app) \\
\midrule
Contacts & Dana's no., Omar's no., \textsc{lm} set (agg.) & edit no., delete, insert \\
SMS & newest from X (agg.), \#unread (agg.), tracked body & insert, edit body, delete newest \\
Call log & newest caller (agg.), \#missed (agg.) & insert, change type, delete \\
MediaStore & newest photo (agg.), tracked names & insert, rename, delete \\
Calendar & stand-up time, \textsc{lm} events (agg.) & move, delete, insert \\
Settings & two \texttt{Global} keys & update \\
Files & \texttt{todo.md} hash, directory list (agg.) & append, create, delete \\
\bottomrule
\end{tabular}
\end{table}

\paragraph{Gap protocol (E2).} A trial ensures base data and waits for deferred notifications from setup to drain. The client then observes all facts, a gap is injected and verified, and each source is mutated with probability 0.5. The gap is held 0--30\,s, ended, and the client decides at +2\,s (\emph{early}) and +14\,s (\emph{late}). Changes are thus 2--32\,s old at the early decision, on both sides of the deferral; a short-hold condition (0--4\,s) is in Appendix~\ref{app:extra}. Kill uses \texttt{am kill} for a background client and \texttt{SIGKILL} for an FGS client. Doze is \texttt{force-idle deep} with the battery unplugged; the restricted standby bucket is a separate gap. Process-ending gaps relaunch the client (and restart its FGS) as the agent would. Two longer gaps run about 10~minutes from observation to decision. In the default-debounce freezer gap the client freezes about 10~minutes after observation; the changes are then made while it is frozen and held for 60\,s (E9 holds 10~minutes). The stepped Doze gap mutates in deep idle, steps through a maintenance window and mutates again; its logged state at the second batch was a maintenance window in 4 of 5 trials, so we repeated it with deep idle verified before both batches (\S\ref{sec:eval-validity}). We ran \ETwoBg background and \ETwoFg FGS trials on Android~13 and \ETwoApi on Android~15.

% AUTO-GENERATED by artifact/scripts/render.py -- do not edit
\begin{table*}[t]
\centering
\caption{Cost vs.\ source size (E12; medians, ms). Contacts: \stamp validation with the VERSION-map adapter (unchanged / after one edit), with the timestamp adapter, and a provider refresh of the 3 Contacts facts (incl.\ the aggregate over all ``LM'' contacts). MediaStore: validation (unchanged / after one insert) and refresh of the 2 media facts (incl.\ the newest-photo aggregate). Last column: boundary check to effect.}
\label{tab:scale}
\mtab\scriptsize
\begin{tabular}{rrrrrrrrrr}
\toprule
\mhead Raw contacts & Images & C.\ val. & C.\ val.\,$\Delta$ & C.\ val.\ TS & C.\ refresh & M.\ val. & M.\ val.\,$\Delta$ & M.\ refresh & Check$\to$effect \\
\midrule
48 & 24 & 7.0 & 15.5 & 9.2 & 9.0 & 4.7 & 12.7 & 5.7 & 19 \\
148 & 34 & 7.1 & 14.4 & 9.6 & 9.1 & 4.7 & 11.3 & 5.7 & 20 \\
1048 & 44 & 9.5 & 22.7 & 10.5 & 13.5 & 5.4 & 14.4 & 5.7 & 23 \\
5048 & 54 & 31.7 & 69.1 & 28.2 & 34.4 & 9.0 & 22.7 & 12.2 & 34 \\
5048 & 1064 & 29.6 & 55.2 & 26.0 & 31.0 & 11.2 & 24.6 & 13.0 & 35 \\
5048 & 10074 & 30.1 & 53.8 & 24.2 & 28.5 & 32.9 & 64.6 & 41.3 & 34 \\
\bottomrule
\end{tabular}
\end{table*}


\paragraph{At scale.} E12 grew the Contacts database to about 5{,}000 raw contacts and the MediaStore folder to about 10{,}000 images on Android~13 (Table~\ref{tab:scale}). Validation reads the whole RawContacts version map, so its cost grows linearly: from \ScaleContactsZeroVc\,ms with \ScaleContactsZeroRaw raw contacts to \ScaleContactsFivezerozerozeroVc\,ms with \ScaleContactsFivezerozerozeroRaw. Refreshing the three Contacts facts grew similarly (\ScaleContactsZeroRc to \ScaleContactsFivezerozerozeroRc\,ms), because one of them is an aggregate over all ``LM'' contacts, which the bulk contacts join. MediaStore validation reads the volume generation and, after a change, a delta; it rose from \ScaleContactsZeroVm to \ScaleMediaOnezerozerozerozeroVm\,ms at \ScaleMediaOnezerozerozerozeroImg images, and refresh from \ScaleContactsZeroRm to \ScaleMediaOnezerozerozerozeroRm\,ms. Validation and refresh stay of the same order at every size, so the conclusion above does not change with scale. The boundary check-to-effect interval grew from \ScaleContactsZeroCte to \ScaleContactsFivezerozerozeroCte\,ms, because the check reads the version map.

\paragraph{The FGS alternative.} An FGS client is the strongest notification-based alternative: it removes the deferral and the freezer (\S\ref{sec:meas-deploy}), and restart-aware event-only on it missed \FgAnyEvrK of \FgAnyChanged changed uses. E11 alternated 4-minute windows with the client cached and with its FGS running, while another app changed a source every 20\,s (Table~\ref{tab:fgcost}). Per window the FGS client used \CostFgCpu\,s of CPU (\CostBgCpu\,s in the background) and \CostFgPss\,MB of PSS (\CostBgPss\,MB); BatteryStats' model put its charge at about \CostFgPower\,mAh per window, and did not list the background client at all. These emulator CPU and BatteryStats estimates do not establish physical energy cost. Deployment also incurs notification and policy constraints. It shows a persistent notification. An app targeting Android~14 must declare an FGS type~\cite{android_fgstypes}; a service must justify an appropriate type; these experiments do not establish eligibility for an always-on memory service. For an app targeting Android~15, \texttt{dataSync} services may run only 6~hours in 24 unless the user opens the app~\cite{android_fgstimeout}. Our client targets API~33, so these rules did not apply to it.

% AUTO-GENERATED by artifact/scripts/render.py -- do not edit
\begin{table*}[t]
\centering
\caption{Cost of keeping the memory in a foreground service (E11, Android~13), per window with a change by another app every 20\,s: medians of CPU time, PSS and BatteryStats' estimated charge for the client's UID (a power-profile estimate on an emulator, not a battery measurement; ``not listed'': below BatteryStats' reporting threshold; in parentheses: windows in which it was listed).}
\label{tab:fgcost}
\mtab
\begin{tabular}{lrrrrr}
\toprule
\mhead Client & Windows & Window (min) & CPU (s) & PSS (MB) & Est. charge (mAh) \\
\midrule
Background (cached) & 3 & 4 & 0.28 & 24.4 & not listed \\
Foreground service & 3 & 4 & 0.21 & 28.8 & $2.6\times10^{-6}$ (2/3) \\
\bottomrule
\end{tabular}
\end{table*}


% AUTO-GENERATED by artifact/scripts/render.py; do not edit
\begin{table*}[t]
\centering
\caption{A host agent reading the phone over adb (E13; Android~13): time per decision (s; median, bootstrap 95\% CI) to re-read all 16 base facts (global refresh, which for base facts is also input memoization) and to check watermarks and re-read what they do not certify (\stamp). Missed: changed facts not re-read.}
\label{tab:host}
\mtab
\begin{tabular}{lrrrrrr}
\toprule
\mhead Decisions & $n$ & Global (s) & \stamp (s) & \stamp faster & Re-reads & Missed \\
\midrule
No change & 20 & 11.9\,[11.9,\,12.0] & 8.0\,[8.0,\,8.0] & 20/20 & 7.0 & 0 \\
One change & 20 & 11.9\,[11.8,\,12.0] & 8.1\,[8.0,\,9.5] & 20/20 & 8.2 & 0 \\
\bottomrule
\end{tabular}
\end{table*}


\paragraph{StampCheck host comparison: unbatched reads.} The stamp-based host agent used no notification relay and issued a separate adb command per read. E13 times both strategies from the host over 40 decisions on Android~13 (Table~\ref{tab:host}). Re-reading all 16 base facts took \HostNoneGl\,s per decision; for base facts this is also what input memoization and restart-aware refresh do. Reading four watermarks and re-reading only what they do not certify (always the 7 SMS, calendar and Settings facts) took \HostNoneEv\,s without a change and \HostChgEv\,s after one, faster in every decision, and missed none of the \HostChgChanged changed facts. The cost is dominated by starting the \texttt{content} tool for each query (median \HostQueryMs\,ms vs.\ \HostShellMs\,ms for a shell command), so an agent that batches its reads into one on-device helper would see a smaller gap. This result establishes an advantage only against the measured unbatched baseline, under its staged mutation protocol. The matched helper comparison below reassesses that advantage.



\subsection{Validity notes}
\label{sec:eval-validity}

Trials are excluded only by outcome-independent rules: harness errors, freezer trials whose cgroup was not frozen, trials whose recorded client mode does not match the condition, and spans longer than the gap allows (host sleep); counts per condition are in the artifact. The shortened freezer debounce must be set with DeviceConfig sync disabled, or sync silently restores the default; freezer trials in the short-hold condition that were not frozen are excluded for this reason. Nine background no-gap trials ran with a foreground service left over from a smoke test and were discarded before analysis. The stepped long-Doze gap was logged as a maintenance window, not deep idle, at the second mutation batch in 4 of 5 trials; the verified repeat steps the idle state until it reads \texttt{IDLE} before each batch and also logs that the client is not frozen (\NLongdozeTwo trials; median first callback \FclLongdozeTwo\,s, max \FclMaxLongdozeTwo\,s). A first run of it is excluded: the freezer debounce was still set to 2\,s from E9, so the client was frozen during the gap and the run measured the freezer. A first run of the persistent app is excluded because its observers were registered before its permissions were granted. Both are kept in the artifact. Our Contacts database came back empty after \RebootEmpty of \RebootN reboots. On Android~15 the 12\,s setup drain of E2 was shorter than the deferral, so some setup notifications arrived after observation and set event flags; this can only make the event-driven policies refresh more, which biases the comparison in their favour (E9 drains for 16\,s).


\subsection{Extended Frozen-Hold and Source Measurements}
E2's freezer gaps are short, so E9 holds the client frozen for 10~minutes with changes pending. Per trial, the client observes and goes to the background; the driver waits until its cgroup is frozen, then another app inserts or updates a random source every 3\,s until 600\,s have passed since the first change (about \NineMutMedA changes per trial). The frozen state is sampled every 10 changes. The driver then sends \textsf{DECIDE}, which thaws the client, and decisions follow at +2, +14 and +120\,s. We ran \NineRunA trials on Android~13 and \NineRunB on Android~15 (Table~\ref{tab:frozen}).

\emph{Observed source invalidations survive these holds.} In the \NineStayedA (Android~13) and \NineStayedB (Android~15) trials that stayed frozen, every one of the \NineSrcMutA and \NineSrcMutB mutated sources got a callback after its last change. Held callbacks were delivered as soon as the process thawed, even before it handled the \textsf{DECIDE} broadcast that thawed it (the first callback came a median of \NineFirstAfterA\,s and \NineFirstAfterB\,s before the client logged that broadcast). Changes made in the last seconds before the thaw were still in ContentService's deferral queue and arrived up to \NineLastAfterMaxA\,s (Android~13) and \NineLastAfterMaxB\,s (Android~15) after it. So restart-aware event-only missed \NineEvrEarlyA (Android~13) and \NineEvrEarlyB (Android~15) changed uses at +2\,s, and none at +14\,s (\NineEvrLateA and \NineEvrLateB). The Android~15 traces also permit re-freezing after thaw. The source-level observations do not establish delivery of every individual callback.

\emph{A long freeze can end in death.} In \NineDiedA of the Android~13 trials (\NineDiedB on Android~15), ActivityManager killed the frozen process \NineDiedAfterA\,s into the hold, with exit reason ``\NineDiedDescA'': the cached-process limit, not the binder buffer. From then on the gap is a \emph{lost} gap: event-only missed \NineDiedEvMissA changed uses even at +120\,s, while restart-aware and \stamp missed none. A memory cannot assume that a freeze, however benign, ends in a thaw.

% AUTO-GENERATED by artifact/scripts/render.py; do not edit
\begin{table*}[t]
\centering
\caption{Historical app/name-based AndroidWorld census. Store classes describe candidate evidence, not certified query coverage. Right column: name-rule candidates for reading pre-existing state; the nine-validator inspection qualifies this rule.}
\label{tab:aw}
\mtab
\begin{tabular}{lrr}
\toprule
\mhead Stores behind the task & All tasks & Reads existing state \\
\midrule
Only stores with durable evidence & 19 & 12 \\
Some API-readable store without evidence & 23 & 6 \\
Some app-private or UI-only store & 74 & 53 \\
\midrule Total & 116 & 71 \\
Touching at least one store with evidence & 28 & 19 \\
\bottomrule
\end{tabular}
\end{table*}

